\section*{Chapter \ref{ch:g11:externalities} --- Externalities}
% figdata: figdata/grade-11/06-externalities/out/ (so2_peaks, so2_usa); Brenland numbers invented.

\begin{solution}{exo:g11:externalities:1}
15, 17, 19 and 21 dalens.
\end{solution}

\begin{solution}{exo:g11:externalities:2}
$300 + 45 = 345$ dalens; the airline bears $300/345 \approx 87\,\%$ of it.
\end{solution}

\begin{solution}{exo:g11:externalities:3}
The United States, in 1918, at about 191 kilograms a person. The United States peaked at about \usd{10500} of GDP per person, Britain at about \usd{7900} and Japan at about
\usd{14100}.
\end{solution}

\begin{solution}{exo:g11:externalities:4}
Market: $80 - Q = Q$, $Q = 40$. Efficient: $80 - Q = Q + 20$, $Q = 30$.
\end{solution}

\begin{solution}{exo:g11:externalities:5}
A tax of 20 a unit. At $Q = 30$ buyers pay $80 - 30 = 50$ and sellers receive $50 - 20 = 30$, their \omterm{def:g11:marginal-reasoning:marginal-benefit}{marginal cost}.
\end{solution}

\begin{solution}{exo:g11:externalities:6}
$\tfrac12 \times (40 - 30) \times 20 = 100$ dalens.
\end{solution}

\begin{solution}{exo:g11:externalities:7}
Market: $130 - 2Q = 10 + Q$, $Q = 40$. Efficient: $130 - 2Q = 10 + 2Q$, $Q = 30$. Tax: the marginal damage at 30, that is 30 a unit. \omterm{def:g10:what-prices-do:dwl}{Deadweight loss}: between 30 and 40
the social \omterm{def:g11:marginal-reasoning:marginal-benefit}{marginal cost} exceeds the demand price by $4Q - 120$, from 0 to 40, so the loss is the triangle $\tfrac12 \times 10 \times 40 = 200$.
\end{solution}

\begin{solution}{exo:g11:externalities:8}
The cheapest outcome is a fence: keeping the goats fenced costs 200, against a loss of 500 if they roam and of 300 if they are given up. If the neighbour has the right,
the farmer pays him between 200 and 500 to build the fence. If the farmer has the right, the neighbour fences his goats himself, since 200 is less than the 300 they are
worth to him. The goats are fenced either way; only who pays changes.
\end{solution}

\begin{solution}{exo:g11:externalities:9}
Now letting the goats roam is cheapest: the damage, 150, is less than the fence, 200. If the neighbour has the right, the farmer would pay at most 150 for a fence costing
200: the goats roam. If the farmer has the right, the neighbour pays him between 150 and 200 a year to let the goats roam, which is cheaper than the fence. The goats roam
either way.
\end{solution}

\begin{solution}{exo:g11:externalities:10}
$x_A = 4x_B$ and $x_A + x_B = 50$ give $x_B = 10$ and $x_A = 40$; the permit price is 40.
\end{solution}

\begin{solution}{exo:g11:externalities:11}
Without a subsidy: $100 - Q/10 = 40$, $Q = 600$. Efficient: $120 - Q/10 = 40$, $Q = 800$. A subsidy of 20 a vaccination, the external benefit, brings the number to 800.
\end{solution}

\begin{solution}{exo:g11:externalities:12}
A case in which rights are clear and the parties are few, yet no deal is made, or an inefficient use continues. Farnsworth's twenty nuisance cases: the court made the
rights clear, but the parties, angry with each other, never bargained; the assumption that bargaining is cheap fails, since resentment is a cost. The Danish waterworks: the
farmers knew their own losses and asked for more; the assumption that both sides know the costs fails.
\end{solution}

\begin{solution}{exo:g11:externalities:13}
Firm 1 cuts $x$ and firm 2 $E - x$; the total cost $C_1(x) + C_2(E - x)$ has derivative $C_1'(x) - C_2'(E - x)$, which is zero where the \omterm{def:g11:marginal-reasoning:marginal-benefit}{marginal costs} are equal, and since
both rise, that is the minimum. Under a tax $t$, each firm cuts while cutting a tonne costs less than the tax it saves, so it stops where $C_i' = t$: both \omterm{def:g11:marginal-reasoning:marginal-benefit}{marginal costs} equal
$t$, and the authority chooses $t$ so that the cuts add up to $E$. A uniform standard gives $E/2$ each, which equalises \omterm{def:g11:marginal-reasoning:marginal-benefit}{marginal costs} only if the two firms have the same costs.
\end{solution}

\begin{solution}{exo:g11:externalities:14}
$100 - Q = 60 + Q$ gives $Q = 20$. Between 20 and 30, the units forgone were worth more than their \omterm{def:g11:externalities:private-cost}{social cost}, by $60 - 2Q$, from 20 to 0: a loss of $\tfrac12 \times 10 \times 20
= 100$, as large as with no tax at all.
\end{solution}

\begin{solution}{exo:g11:externalities:15}
The external benefit of a vaccination is the infections it prevents in others. When the disease spreads fast and few are protected, one vaccination prevents many infections,
worth far more than 20. When nearly everyone is vaccinated, the disease can hardly spread, and one more vaccination protects almost no one else: the external benefit falls
towards zero. A fixed subsidy is too small in the first case and too large in the second.
\end{solution}

\begin{solution}{pb:g11:externalities:1}
\textbf{1.} $120 - 2q = 0$: $q = 60$ tonnes.

\textbf{2.} $120 - 2q = q$: $q = 40$ tonnes.

\textbf{3.} The damage is the area under $q$: $q^2/2$, that is 1800 at 60 and 800 at 40.

\textbf{4.} Between 40 and 60 the marginal damage exceeds the \omterm{def:g11:marginal-reasoning:marginal-benefit}{marginal benefit} by $3q - 120$, from 0 to 60: the loss is $\tfrac12 \times 20 \times 60 = 600$.

\textbf{5.} $B(q) = 120q - q^2$: 3600 at 60 and 3200 at 40.

\textbf{6.} The marginal damage at the efficient quantity: 40 dalens a tonne.

\textbf{7.} $120 - 2q = 40$ gives $q = 40$.

\textbf{8.} $40 \times 40 = 1600$ dalens a week.

\textbf{9.} All of it: the market reaches the efficient quantity, and the 600 dalens are recovered.

\textbf{10.} It would set the quantity where $120 - 2q = 2q$, $q = 30$, with a tax of 60. With the true damage, that cuts too far: the loss between 30 and 40 is $\tfrac12 \times 10
\times 30 = 150$.

\textbf{11.} 40 tonnes. The fishery accepts if it is paid at least its damage, 800; the tannery gains 3200 from 40 tonnes and pays.

\textbf{12.} 40 tonnes. Cutting from 60 to 40 costs the tannery $3600 - 3200 = 400$ and saves the fishery $1800 - 800 = 1000$: the fishery pays between 400 and 1000.

\textbf{13.} The quantity is 40 in both. Who pays changes: with the fishery's right the tannery pays at least 800; with the tannery's right the fishery pays between 400 and 1000.

\textbf{14.} A thousand households would have to agree among themselves and with the tannery; each would hope the others pay, and the cost of reaching and enforcing a deal
would exceed the gain.

\textbf{15.} For example: measuring the waste and the harm to the fish; lawyers and the time spent negotiating; checking afterwards that the tannery keeps to the deal.

\textbf{16.} A: $\int_0^{20} x\,dx = 200$; B: $\int_0^{20} 4x\,dx = 800$; total 1000.

\textbf{17.} $x_A = 4x_B$ and $x_A + x_B = 40$ give $x_B = 8$, $x_A = 32$; the price is 32 dalens a tonne.

\textbf{18.} A: $32^2/2 = 512$; B: $2 \times 8^2 = 128$; total 640.

\textbf{19.} $1000 - 640 = 360$ dalens a week, 36\,\%.

\textbf{20.} The efficient tax is 40 dalens a tonne, and it recovers the whole \omterm{def:g10:what-prices-do:dwl}{deadweight loss} of 600 dalens a week; permits for the same total cut cost 360 dalens a week (36\,\%) less
than a uniform standard.
\end{solution}
